\documentclass[UTF8,a4paper,11pt]{ctexart}
\usepackage[a4paper,top=2.35cm,bottom=2.1cm,left=2.3cm,right=2.3cm]{geometry}
\usepackage{amsmath,amssymb,mathtools,booktabs,array,xcolor}
\usepackage{fancyhdr,lastpage,hyperref}
\definecolor{ink}{HTML}{15324B}
\definecolor{accent}{HTML}{167D8D}
\definecolor{dim}{HTML}{506778}
\hypersetup{colorlinks=true,linkcolor=accent,urlcolor=accent,pdfauthor={三千弱水大学物理原创试题库}}
\pagestyle{fancy}\fancyhf{}\fancyhead[L]{\small\color{dim} 三千弱水 · 大学物理原创题库}
\fancyhead[R]{\small\color{dim} 阶段一 · 样板题}\fancyfoot[C]{\small\thepage/\pageref{LastPage}}
\renewcommand{\headrulewidth}{0.3pt}
\setlength{\parindent}{2em}\setlength{\parskip}{.35em}
\allowdisplaybreaks
\hypersetup{pdftitle={大学物理原创试题库 · 参考解答}}
\begin{document}
\begin{titlepage}
\thispagestyle{empty}
\vspace*{1.5cm}
{\color{accent}\rule{\linewidth}{2.4pt}}
\vspace{2.1cm}
\begin{center}
{\fontsize{32}{42}\selectfont\bfseries\color{ink} 大学物理原创试题库}\par
\vspace{.55cm}
{\fontsize{21}{30}\selectfont\color{ink} 参考解答与推导}\par
\vspace{.35cm}
{\large\color{dim} Solutions and Derivations}\par
\vspace{2.0cm}
{\Huge\color{accent}\bfseries 36\quad 道}\par
\vspace{.65cm}
{\large 力学 · 电磁学 · 热学 · 振动与波动 · 光学 · 近代物理}\par
\vspace{1.25cm}
{\color{dim}\rule{.55\linewidth}{.6pt}}\par
\vspace{.8cm}
{\large 三千弱水 · Ruoshui Physics Bank}\par
\vspace{.3cm}
{\color{dim}\small 2026 年 10 月 · 第一批样板 · 不代表六轮人工质检已完成}\par
\end{center}
\vfill
{\color{accent}\rule{\linewidth}{2.4pt}}
\end{titlepage}
\tableofcontents\clearpage
\section{力学}
\noindent{\large\bfseries\color{ink} M001 · 加速电梯中的落体}\quad{\small\color{dim}[难度 2 / 5 · 计算题]}\par
\noindent\textbf{题目}：一电梯以恒定加速度 $a>0$ 竖直向上运动。相对电梯静止的小球从距地板 $h$ 处无初速释放。忽略空气阻力，求小球落到地板的时间。
\par\noindent\textbf{解答}：取电梯为非惯性参考系，竖直向下为正。小球相对电梯受到的等效加速度为 $g+a$。初始相对速度为零，故 $h=\frac12(g+a)t^2$，由 $h>0$、$g+a>0$ 得唯一正根。
\par\noindent\textbf{结论}：\(\displaystyle t=\sqrt{\frac{2h}{g+a}}\)\par
\vspace{1.0em}\par\noindent{\color{accent}\rule{\linewidth}{.35pt}}\vspace{.8em}
\noindent{\large\bfseries\color{ink} M002 · 传送带上的滑移}\quad{\small\color{dim}[难度 2 / 5 · 计算题]}\par
\noindent\textbf{题目}：水平传送带以恒速 $v>0$ 运动，质量为 $m$ 的物块轻放在带上，初速为零。动摩擦因数为 $\mu>0$，求达到带速所需时间和其间物块相对传送带的滑移距离。
\par\noindent\textbf{解答}：滑动阶段物块受恒定摩擦力 $\mu mg$，加速度 $\mu g$。由 $v=\mu gt$ 得 $t=v/(\mu g)$。这段时间带面位移 $vt$，物块位移 $\tfrac12\mu gt^2$，两者差即滑移距离 $vt-\tfrac12\mu gt^2$。
\par\noindent\textbf{结论}：\(\displaystyle t=\frac{v}{\mu g},\quad s_{\mathrm{rel}}=\frac{v^2}{2\mu g}\)\par
\vspace{1.0em}\par\noindent{\color{accent}\rule{\linewidth}{.35pt}}\vspace{.8em}
\noindent{\large\bfseries\color{ink} M003 · 圆柱纯滚动下坡}\quad{\small\color{dim}[难度 3 / 5 · 推导题]}\par
\noindent\textbf{题目}：质量为 $m$、半径为 $R$ 的均匀实心圆柱，在倾角 $\theta$ 的粗糙斜面上由静止无滑动滚下，求质心加速度；下降竖直高度 $H$ 后的速度。设静摩擦足够大。
\par\noindent\textbf{解答}：沿斜面有 $mg\sin\theta-f=ma$，绕中心有 $fR=I\alpha$，且 $I=\tfrac12mR^2$、$a=R\alpha$。联立解得 $a=\tfrac23g\sin\theta$。机械能守恒：$mgH=\tfrac12mv^2+\tfrac12I(v/R)^2=\tfrac34mv^2$。
\par\noindent\textbf{结论}：\(\displaystyle a=\frac23g\sin\theta,\quad v=\sqrt{\frac{4gH}{3}}\)\par
\vspace{1.0em}\par\noindent{\color{accent}\rule{\linewidth}{.35pt}}\vspace{.8em}
\noindent{\large\bfseries\color{ink} M004 · 完全非弹性碰撞能损}\quad{\small\color{dim}[难度 2 / 5 · 计算题]}\par
\noindent\textbf{题目}：光滑水平面上，质量为 $m$ 的小车以速度 $v$ 撞上静止的质量为 $2m$ 的小车，碰后两车粘在一起。求共同速度及损失的机械能。
\par\noindent\textbf{解答}：水平方向动量守恒，$mv=(m+2m)u$，故 $u=v/3$。碰前动能 $K_i=mv^2/2$，碰后 $K_f=(3m)(v/3)^2/2=mv^2/6$，损失 $K_i-K_f=mv^2/3$。
\par\noindent\textbf{结论}：\(\displaystyle u=\frac v3,\quad \Delta E_{\mathrm{loss}}=\frac{mv^2}{3}\)\par
\vspace{1.0em}\par\noindent{\color{accent}\rule{\linewidth}{.35pt}}\vspace{.8em}
\noindent{\large\bfseries\color{ink} M005 · 位置依赖的驱动力与末速}\quad{\small\color{dim}[难度 3 / 5 · 计算题]}\par
\noindent\textbf{题目}：质量为 $m>0$ 的质点在水平光滑直轨道上运动，受沿正 $x$ 方向的力 $F(x)=kx$（$k>0$）。质点从 $x=a>0$ 静止出发，运动至 $x=b>a$。求到达 $b$ 时的速率和通过该区间的功。
\par\noindent\textbf{解答}：因为轨道光滑，合外力所做的功全部转化为动能。由积分 $W=\int_a^b kx\,\mathrm dx=\frac{k}{2}(b^2-a^2)$；初动能为零，因此 $\frac12mv^2=W$。取正平方根，运动方向始终为正 $x$。
\par\noindent\textbf{结论}：\(\displaystyle W=\frac{k}{2}(b^2-a^2),\quad v=\sqrt{\frac{k(b^2-a^2)}m}\)\par
\vspace{1.0em}\par\noindent{\color{accent}\rule{\linewidth}{.35pt}}\vspace{.8em}
\noindent{\large\bfseries\color{ink} M006 · 计入滑轮转动惯量的阿特伍德机}\quad{\small\color{dim}[难度 4 / 5 · 推导题]}\par
\noindent\textbf{题目}：两物体质量分别为 $m_1,m_2$，且 $m_2>m_1>0$，由质量可忽略且不可伸长的细绳跨过半径为 $R$、转动惯量为 $I$ 的滑轮连接。绳不打滑，滑轮轴无摩擦，求加速度及两侧张力。
\par\noindent\textbf{解答}：设 $m_2$ 向下、$m_1$ 向上运动，大小加速度均为 $a$。两物体满足 $T_1-m_1g=m_1a$、$m_2g-T_2=m_2a$。滑轮满足 $(T_2-T_1)R=I\alpha$，无滑动给出 $\alpha=a/R$。将 $T_1=m_1(g+a)$、$T_2=m_2(g-a)$ 代入转动方程，得 $(m_2-m_1)g=(m_1+m_2+I/R^2)a$。
\par\noindent\textbf{结论}：\(\displaystyle a=\frac{(m_2-m_1)g}{m_1+m_2+I/R^2},\quad T_1=m_1(g+a),\quad T_2=m_2(g-a)\)\par
\vspace{1.0em}\par\noindent{\color{accent}\rule{\linewidth}{.35pt}}\vspace{.8em}
\section{电磁学}
\noindent{\large\bfseries\color{ink} E001 · 非均匀带电细杆的轴向场}\quad{\small\color{dim}[难度 4 / 5 · 推导题]}\par
\noindent\textbf{题目}：细杆沿 $x$ 轴从 $x=0$ 延至 $L$，线电荷密度 $\lambda(x)=\lambda_0 x/L$，其中 $\lambda_0>0$。求 $x=2L$ 处的电场大小及方向。
\par\noindent\textbf{解答}：各电荷元在观察点产生的电场均沿 $+x$。由库仑定律 $E=\frac{1}{4\pi\varepsilon_0}\frac{\lambda_0}{L}\int_0^L\frac{x}{(2L-x)^2}\,\mathrm dx$。令 $u=2L-x$，积分化为 $\left[\frac{2L}{2L-x}+\ln(2L-x)\right]_0^L=1-\ln2$；对数在上下限之差中无量纲化。
\par\noindent\textbf{结论}：\(\displaystyle E=\frac{\lambda_0(1-\ln2)}{4\pi\varepsilon_0 L}\quad(\text{方向 }+x)\)\par
\vspace{1.0em}\par\noindent{\color{accent}\rule{\linewidth}{.35pt}}\vspace{.8em}
\noindent{\large\bfseries\color{ink} E002 · 无限长均匀带电圆柱}\quad{\small\color{dim}[难度 3 / 5 · 推导题]}\par
\noindent\textbf{题目}：半径 $R$ 的无限长非导体实心圆柱均匀带正体电荷密度 $\rho$。忽略端部效应，分别求柱内 $r<R$、柱外 $r>R$ 的电场。
\par\noindent\textbf{解答}：取长度为 $\ell$、半径为 $r$ 的同轴高斯面。由对称性 $E(2\pi r\ell)=Q_{\mathrm{enc}}/\varepsilon_0$。柱内 $Q_{\mathrm{enc}}=\rho\pi r^2\ell$；柱外 $Q_{\mathrm{enc}}=\rho\pi R^2\ell$。方向均沿径向向外，两式在 $r=R$ 处连续。
\par\noindent\textbf{结论}：\(\displaystyle E(r)=\begin{cases}\rho r/(2\varepsilon_0),&r\le R,\\ \rho R^2/(2\varepsilon_0r),&r\ge R.\end{cases}\)\par
\vspace{1.0em}\par\noindent{\color{accent}\rule{\linewidth}{.35pt}}\vspace{.8em}
\noindent{\large\bfseries\color{ink} E003 · 断开电源后的介质插入}\quad{\small\color{dim}[难度 2 / 5 · 概念推导题]}\par
\noindent\textbf{题目}：一平行板电容器先充电，随后与电源断开。忽略边缘效应，将相对介电常数为 $\kappa>1$ 的均匀电介质完全填入极板间。求电容、电压及储能相对原值的比。
\par\noindent\textbf{解答}：断开电源后自由电荷量 $Q$ 不变，填充前后电容分别为 $C_0$ 与 $\kappa C_0$。利用 $V=Q/C$ 和 $U=Q^2/(2C)$，可得电压与储能均变为原来的 $1/\kappa$。
\par\noindent\textbf{结论}：\(\displaystyle C/C_0=\kappa,\quad V/V_0=U/U_0=1/\kappa\)\par
\vspace{1.0em}\par\noindent{\color{accent}\rule{\linewidth}{.35pt}}\vspace{.8em}
\noindent{\large\bfseries\color{ink} E004 · 匀强磁场中的旋转线圈}\quad{\small\color{dim}[难度 2 / 5 · 计算题]}\par
\noindent\textbf{题目}：单匝面积为 $A$、总电阻为 $R$ 的矩形闭合线圈，以角速度 $\omega$ 绕线圈平面内的一条轴匀速转动。匀强磁场大小为 $B$。在 $t=0$ 时磁场与线圈法向平行，求感应电动势与电流。忽略自感。
\par\noindent\textbf{解答}：选取面积法向的初始方向，磁通量 $\Phi(t)=BA\cos(\omega t)$。法拉第定律给出 $\mathcal E=-\mathrm d\Phi/\mathrm dt=BA\omega\sin(\omega t)$，由欧姆定律 $I=\mathcal E/R$。正负号表示相对所选绕向的电流方向。
\par\noindent\textbf{结论}：\(\displaystyle \mathcal E=BA\omega\sin(\omega t),\quad I=\frac{BA\omega}{R}\sin(\omega t)\)\par
\vspace{1.0em}\par\noindent{\color{accent}\rule{\linewidth}{.35pt}}\vspace{.8em}
\noindent{\large\bfseries\color{ink} E005 · 同轴圆柱电容器的单位长度电容}\quad{\small\color{dim}[难度 3 / 5 · 推导题]}\par
\noindent\textbf{题目}：两同轴理想导体圆柱半径为 $a$、$b$（$0<a<b$），长度 $L\gg b$，其间充满介电常数为 $\varepsilon$ 的线性均匀介质。内外导体分别带单位长度电荷 $+\lambda$、$-\lambda$，忽略端部效应，求单位长度电容及电场能量。
\par\noindent\textbf{解答}：对 $a<r<b$ 作同轴高斯圆柱面，由高斯定律 $E_r\,2\pi r\ell=\lambda\ell/\varepsilon$ 得 $E_r=\lambda/(2\pi\varepsilon r)$。电势差 $V=\int_a^b E_r\,\mathrm dr=\lambda\ln(b/a)/(2\pi\varepsilon)$。定义 $C/L=\lambda/V$。单位长度储能 $U/L=\tfrac12\lambda V$。
\par\noindent\textbf{结论}：\(\displaystyle \frac CL=\frac{2\pi\varepsilon}{\ln(b/a)},\quad\frac UL=\frac{\lambda^2}{4\pi\varepsilon}\ln\frac ba\)\par
\vspace{1.0em}\par\noindent{\color{accent}\rule{\linewidth}{.35pt}}\vspace{.8em}
\noindent{\large\bfseries\color{ink} E006 · 时变圆形磁场激发的环形电场}\quad{\small\color{dim}[难度 4 / 5 · 推导题]}\par
\noindent\textbf{题目}：在理想化模型中，沿 $z$ 方向的均匀磁场仅存在于半径 $R$ 的圆柱区域 $r<R$ 内，其磁感应强度按 $B(t)=B_0+\beta t$ 变化，圆柱外磁场为零。忽略端部效应。求空间中感应电场的环向分量 $E_\varphi(r)$，规定正方向为从 $+z$ 方向观察的逆时针方向。
\par\noindent\textbf{解答}：以半径 $r$ 的同心圆为积分回路。法拉第定律给出 $2\pi rE_\varphi=-\mathrm d\Phi_B/\mathrm dt$。若 $r\le R$，回路所围磁通 $\Phi_B=\pi r^2B(t)$；若 $r\ge R$，磁通 $\Phi_B=\pi R^2B(t)$。分别代入，负号表示当 $\beta>0$ 时沿顺时针方向。
\par\noindent\textbf{结论}：\(\displaystyle E_\varphi(r)=\begin{cases}-\beta r/2,&r\le R,\\-\beta R^2/(2r),&r\ge R.\end{cases}\)\par
\vspace{1.0em}\par\noindent{\color{accent}\rule{\linewidth}{.35pt}}\vspace{.8em}
\section{热学}
\noindent{\large\bfseries\color{ink} T001 · 绝热系统中的热平衡}\quad{\small\color{dim}[难度 3 / 5 · 推导题]}\par
\noindent\textbf{题目}：一绝热容器由固定、导热但不透气的隔板分成两室，内装同一种单原子理想气体，物质的量为 $n_1,n_2$，初温为 $T_1,T_2$。求最终共同温度。两室体积保持不变。
\par\noindent\textbf{解答}：整个系统与外界无热量及功交换，且两室气体均不发生体积功。单原子理想气体内能 $U=\tfrac32nRT$，所以 $\tfrac32n_1RT_1+\tfrac32n_2RT_2=\tfrac32(n_1+n_2)RT_f$。因为隔板不透气，过程中无混合熵项。
\par\noindent\textbf{结论}：\(\displaystyle T_f=\frac{n_1T_1+n_2T_2}{n_1+n_2}\)\par
\vspace{1.0em}\par\noindent{\color{accent}\rule{\linewidth}{.35pt}}\vspace{.8em}
\noindent{\large\bfseries\color{ink} T002 · 卡诺热机热量核算}\quad{\small\color{dim}[难度 2 / 5 · 计算题]}\par
\noindent\textbf{题目}：理想卡诺热机在 $T_h=500\,\mathrm K$ 和 $T_c=300\,\mathrm K$ 的两个热源间工作，每循环输出净功 $400\,\mathrm J$。求效率、从高温热源吸热及向低温热源放热。
\par\noindent\textbf{解答}：卡诺效率 $\eta=1-T_c/T_h=1-300/500=0.4$。由 $W=\eta Q_h$ 得 $Q_h=400/0.4=1000\,\mathrm J$。热机循环内能变化为零，故 $Q_c=Q_h-W=600\,\mathrm J$。
\par\noindent\textbf{结论}：\(\displaystyle \eta=40\%,\quad Q_h=1000\,\mathrm J,\quad Q_c=600\,\mathrm J\)\par
\vspace{1.0em}\par\noindent{\color{accent}\rule{\linewidth}{.35pt}}\vspace{.8em}
\noindent{\large\bfseries\color{ink} T003 · 单原子气体的等压升温}\quad{\small\color{dim}[难度 2 / 5 · 推导题]}\par
\noindent\textbf{题目}：$n$ mol 单原子理想气体在压强不变时，由温度 $T_0$ 膨胀至体积为原来的两倍。求气体对外做功 $W$、内能增量 $\Delta U$ 和吸热量 $Q$。
\par\noindent\textbf{解答}：等压下由状态方程 $V\propto T$，故 $T_f=2T_0$、$\Delta T=T_0$。于是 $W=p\Delta V=nR\Delta T=nRT_0$，$\Delta U=nC_V\Delta T=\tfrac32nRT_0$，热力学第一定律给出 $Q=\Delta U+W$。
\par\noindent\textbf{结论}：\(\displaystyle W=nRT_0,\quad \Delta U=\frac32nRT_0,\quad Q=\frac52nRT_0\)\par
\vspace{1.0em}\par\noindent{\color{accent}\rule{\linewidth}{.35pt}}\vspace{.8em}
\noindent{\large\bfseries\color{ink} T004 · 理想气体的自由膨胀}\quad{\small\color{dim}[难度 3 / 5 · 推导题]}\par
\noindent\textbf{题目}：容器绝热且刚性，其中 $n$ mol 理想气体初占体积 $V$，抽去隔板后向真空自由膨胀至 $2V$。求温度变化、内能变化及熵变。
\par\noindent\textbf{解答}：自由膨胀不做外功且绝热，因此 $\Delta U=Q-W=0$。理想气体内能仅依赖温度，所以 $\Delta T=0$。熵变通过连接同始末态的可逆等温过程计算，$\Delta S=\int_V^{2V}nR\,\mathrm dV/V=nR\ln2>0$。
\par\noindent\textbf{结论}：\(\displaystyle \Delta T=0,\quad \Delta U=0,\quad \Delta S=nR\ln2\)\par
\vspace{1.0em}\par\noindent{\color{accent}\rule{\linewidth}{.35pt}}\vspace{.8em}
\noindent{\large\bfseries\color{ink} T005 · 理想气体可逆等温压缩}\quad{\small\color{dim}[难度 2 / 5 · 计算题]}\par
\noindent\textbf{题目}：$n$ 摩尔理想气体在温度 $T>0$ 下准静态等温地从体积 $V_1$ 压缩到 $V_2<V_1$。采用“气体对外做功为正”的约定，求气体做功、吸收的热量及内能变化。
\par\noindent\textbf{解答}：理想气体内能仅是温度函数，故 $\Delta U=0$。可逆路径的压强 $p=nRT/V$，气体做功 $W=\int_{V_1}^{V_2}p\,\mathrm dV=nRT\ln(V_2/V_1)<0$。由热力学第一定律 $\Delta U=Q-W$ 得 $Q=W<0$，即气体向外放热。
\par\noindent\textbf{结论}：\(\displaystyle W=Q=nRT\ln\frac{V_2}{V_1}<0,\quad\Delta U=0\)\par
\vspace{1.0em}\par\noindent{\color{accent}\rule{\linewidth}{.35pt}}\vspace{.8em}
\noindent{\large\bfseries\color{ink} T006 · 等热容量气体的隔热热接触}\quad{\small\color{dim}[难度 4 / 5 · 推导题]}\par
\noindent\textbf{题目}：两个刚性容器分别装有相同物质的量 $n$ 的同种理想气体，定容摩尔热容相同且为常数 $C_V$。初温为 $T_1,T_2>0$。打开两容器之间的导热连接（无气体交换），整套装置与外界绝热。求最终共同温度及体系的总熵变。
\par\noindent\textbf{解答}：两个容器均保持体积不变，且系统不与外界交换热量，能量守恒给出 $nC_V(T_f-T_1)+nC_V(T_f-T_2)=0$，于是 $T_f=(T_1+T_2)/2$。分别用可逆定容路径计算熵变 $\Delta S_i=nC_V\ln(T_f/T_i)$；相加得 $\Delta S=nC_V\ln(T_f^2/(T_1T_2))$，由算术—几何平均不等式知 $\Delta S\ge0$。
\par\noindent\textbf{结论}：\(\displaystyle T_f=\frac{T_1+T_2}{2},\quad\Delta S=nC_V\ln\frac{(T_1+T_2)^2}{4T_1T_2}\ge0\)\par
\vspace{1.0em}\par\noindent{\color{accent}\rule{\linewidth}{.35pt}}\vspace{.8em}
\section{振动与波动}
\noindent{\large\bfseries\color{ink} W001 · 行波的传播方向}\quad{\small\color{dim}[难度 2 / 5 · 计算题]}\par
\noindent\textbf{题目}：一平面简谐波位移为 $y(x,t)=0.020\cos(4\pi x-20\pi t)$（SI制）。求振幅、波长、频率、波速及传播方向。
\par\noindent\textbf{解答}：与标准式 $y=A\cos(kx-\omega t)$ 比较，$A=0.020\,\mathrm m$，$k=4\pi\,\mathrm{rad/m}$，$\omega=20\pi\,\mathrm{rad/s}$。于是 $\lambda=2\pi/k=0.50\,\mathrm m$，$f=\omega/(2\pi)=10\,\mathrm{Hz}$，$v=\omega/k=5.0\,\mathrm{m/s}$。相位 $kx-\omega t$ 对应沿 $+x$ 传播。
\par\noindent\textbf{结论}：\(\displaystyle A=0.020\,\mathrm m,\ \lambda=0.50\,\mathrm m,\ f=10\,\mathrm{Hz},\ v=5.0\,\mathrm{m/s}\)\par
\vspace{1.0em}\par\noindent{\color{accent}\rule{\linewidth}{.35pt}}\vspace{.8em}
\noindent{\large\bfseries\color{ink} W002 · 弹簧连接方式与周期}\quad{\small\color{dim}[难度 2 / 5 · 推导题]}\par
\noindent\textbf{题目}：两根劲度系数均为 $k$ 的轻弹簧，分别以串联和并联方式连接相同的质量 $m$，作水平无阻尼简谐运动。求两种方式的周期之比。
\par\noindent\textbf{解答}：串联时受力相等、总伸长相加，故 $k_s=k/2$；并联时伸长相同、回复力相加，故 $k_p=2k$。周期 $T=2\pi\sqrt{m/k_{\mathrm{eff}}}$，因此 $T_s/T_p=\sqrt{k_p/k_s}=2$。
\par\noindent\textbf{结论}：\(\displaystyle \frac{T_{\mathrm{串}}}{T_{\mathrm{并}}}=2\)\par
\vspace{1.0em}\par\noindent{\color{accent}\rule{\linewidth}{.35pt}}\vspace{.8em}
\noindent{\large\bfseries\color{ink} W003 · 运动声源的多普勒效应}\quad{\small\color{dim}[难度 2 / 5 · 计算题]}\par
\noindent\textbf{题目}：静止观察者面前，一频率为 $500\,\mathrm{Hz}$ 的声源沿连线方向以 $20\,\mathrm{m/s}$ 向观察者匀速运动。声速为 $340\,\mathrm{m/s}$，求观察频率。忽略风。
\par\noindent\textbf{解答}：声源接近静止观察者时，相邻波峰的空间间隔为 $\lambda'=(c-v_s)/f$。观察者测得 $f'=c/\lambda'=cf/(c-v_s)=340\times500/(340-20)$。
\par\noindent\textbf{结论}：\(\displaystyle f'=531.25\,\mathrm{Hz}\)\par
\vspace{1.0em}\par\noindent{\color{accent}\rule{\linewidth}{.35pt}}\vspace{.8em}
\noindent{\large\bfseries\color{ink} W004 · 两端固定弦的第三谐波}\quad{\small\color{dim}[难度 2 / 5 · 概念推导题]}\par
\noindent\textbf{题目}：长度为 $L$ 的均匀细弦两端固定，横波速度为 $v$。求第三谐波频率及弦内部（不含两端）有多少个节点。
\par\noindent\textbf{解答}：两端固定弦允许波长 $\lambda_n=2L/n$，对应 $f_n=nv/(2L)$。第三谐波在 $x=jL/3$（$j=0,1,2,3$）处有节点，排除两个端点后，内部节点为 $x=L/3,2L/3$。
\par\noindent\textbf{结论}：\(\displaystyle f_3=\frac{3v}{2L},\quad \text{内部节点数}=2\)\par
\vspace{1.0em}\par\noindent{\color{accent}\rule{\linewidth}{.35pt}}\vspace{.8em}
\noindent{\large\bfseries\color{ink} W005 · 双频振动的拍频}\quad{\small\color{dim}[难度 2 / 5 · 推导题]}\par
\noindent\textbf{题目}：同一点上叠加两同振幅、同初相、振动方向相同的简谐振动 $x_1=A\cos(2\pi f_1t)$、$x_2=A\cos(2\pi f_2t)$，其中 $f_1,f_2>0$ 且很接近。写出合振动，并求拍频。
\par\noindent\textbf{解答}：利用和差化积：$x=x_1+x_2=2A\cos[\pi(f_1-f_2)t]\cos[\pi(f_1+f_2)t]$。快速振荡的包络振幅为 $2A|\cos[\pi(f_1-f_2)t]|$，相邻振幅极大值的间隔为 $1/|f_1-f_2|$。因此拍频为频率差的绝对值。
\par\noindent\textbf{结论}：\(\displaystyle f_{\mathrm{beat}}=|f_1-f_2|\)\par
\vspace{1.0em}\par\noindent{\color{accent}\rule{\linewidth}{.35pt}}\vspace{.8em}
\noindent{\large\bfseries\color{ink} W006 · 拉紧弦的本征频率}\quad{\small\color{dim}[难度 3 / 5 · 推导题]}\par
\noindent\textbf{题目}：一均匀细弦长为 $L$，两端固定，张力为 $F$，线密度为 $\mu$，小幅横振的波速为 $v=\sqrt{F/\mu}$。求全部非零固有频率、相邻固有频率间距以及第三阶模态的弦内节点个数。
\par\noindent\textbf{解答}：两端位移为零意味着 $L=n\lambda_n/2$（$n=1,2,\dots$），故 $\lambda_n=2L/n$，$f_n=v/\lambda_n=n\sqrt{F/\mu}/(2L)$，频率等间隔。第 $n$ 阶两固定端之外有 $n-1$ 个内部节点，第三阶有2个。
\par\noindent\textbf{结论}：\(\displaystyle f_n=\frac{n}{2L}\sqrt{\frac F\mu},\quad f_{n+1}-f_n=\frac{1}{2L}\sqrt{\frac F\mu},\quad N_{\mathrm{inner},3}=2\)\par
\vspace{1.0em}\par\noindent{\color{accent}\rule{\linewidth}{.35pt}}\vspace{.8em}
\section{光学}
\noindent{\large\bfseries\color{ink} O001 · 两次相位反转的薄膜反射}\quad{\small\color{dim}[难度 3 / 5 · 推导题]}\par
\noindent\textbf{题目}：空气、折射率为 $n_f=1.50$ 的薄膜、折射率为 $n_s=1.70$ 的基底依次叠放。波长为 $600\,\mathrm{nm}$ 的单色光垂直入射，求反射光第一次出现非零膜厚的相长干涉对应的最小膜厚。忽略多次反射。
\par\noindent\textbf{解答}：空气到薄膜、薄膜到基底的反射均为由疏到密，各产生 $\pi$ 相位反转，二者附加相位差抵消。两束反射光光程差为 $2n_ft$，相长条件 $2n_ft=m\lambda$。最小正厚度取 $m=1$，所以 $t=600/(2\times1.50)\,\mathrm{nm}$。
\par\noindent\textbf{结论}：\(\displaystyle t_{\min}=200\,\mathrm{nm}\)\par
\vspace{1.0em}\par\noindent{\color{accent}\rule{\linewidth}{.35pt}}\vspace{.8em}
\noindent{\large\bfseries\color{ink} O002 · 单缝衍射首暗纹}\quad{\small\color{dim}[难度 2 / 5 · 计算题]}\par
\noindent\textbf{题目}：波长 $500\,\mathrm{nm}$ 的单色光垂直照射宽度为 $0.20\,\mathrm{mm}$ 的单缝，观察屏位于缝后 $2.0\,\mathrm m$。在小角近似下求中央亮纹中心到第一级暗纹的距离。
\par\noindent\textbf{解答}：夫琅禾费单缝第一级暗纹满足 $a\sin\theta=\lambda$。小角度时 $y\simeq L\theta\simeq L\lambda/a$。代入 $L=2.0\,\mathrm m$、$\lambda=5.0\times10^{-7}\,\mathrm m$、$a=2.0\times10^{-4}\,\mathrm m$，得 $y=0.0050\,\mathrm m$。
\par\noindent\textbf{结论}：\(\displaystyle y_1\simeq5.0\,\mathrm{mm}\)\par
\vspace{1.0em}\par\noindent{\color{accent}\rule{\linewidth}{.35pt}}\vspace{.8em}
\noindent{\large\bfseries\color{ink} O003 · 杨氏双缝条纹间距}\quad{\small\color{dim}[难度 2 / 5 · 计算题]}\par
\noindent\textbf{题目}：两缝间距为 $d=0.50\,\mathrm{mm}$，屏距 $L=2.0\,\mathrm m$，入射真空波长为 $500\,\mathrm{nm}$ 的相干单色光，且小角近似成立。求相邻明纹间距。
\par\noindent\textbf{解答}：相长条件 $d\sin\theta=m\lambda$。小角近似下相邻级数在屏上的距离差 $\Delta y=L\lambda/d=2.0\times5.0\times10^{-7}/(5.0\times10^{-4})\,\mathrm m$。
\par\noindent\textbf{结论}：\(\displaystyle \Delta y=2.0\,\mathrm{mm}\)\par
\vspace{1.0em}\par\noindent{\color{accent}\rule{\linewidth}{.35pt}}\vspace{.8em}
\noindent{\large\bfseries\color{ink} O004 · 布儒斯特角}\quad{\small\color{dim}[难度 2 / 5 · 计算题]}\par
\noindent\textbf{题目}：光从折射率为 $n_1=1$ 的介质射向折射率为 $n_2=\sqrt3$ 的非吸收介质，求使反射光完全线偏振的入射角（布儒斯特角）。
\par\noindent\textbf{解答}：布儒斯特条件下反射光与折射光互相垂直，由折射定律可得 $\tan\theta_B=n_2/n_1=\sqrt3$，且 $0<\theta_B<\pi/2$。
\par\noindent\textbf{结论}：\(\displaystyle \theta_B=60^\circ\)\par
\vspace{1.0em}\par\noindent{\color{accent}\rule{\linewidth}{.35pt}}\vspace{.8em}
\noindent{\large\bfseries\color{ink} O005 · 凸透镜的物像位置和放大率}\quad{\small\color{dim}[难度 2 / 5 · 计算题]}\par
\noindent\textbf{题目}：薄凸透镜焦距为 $f>0$，一实物垂直光轴放在透镜前方 $s=3f$ 处，采用近轴成像。求像距、像的性质和横向放大率。
\par\noindent\textbf{解答}：薄透镜公式 $1/f=1/s+1/s^{\prime}$，代入 $s=3f$ 得 $1/s^{\prime}=2/(3f)$，故 $s^{\prime}=3f/2>0$，为另一侧实像。横向放大率 $m=-s^{\prime}/s=-1/2$，表示倒立且缩小。
\par\noindent\textbf{结论}：\(\displaystyle s^{\prime}=\frac{3f}{2},\quad m=-\frac12\quad(\text{倒立实像})\)\par
\vspace{1.0em}\par\noindent{\color{accent}\rule{\linewidth}{.35pt}}\vspace{.8em}
\noindent{\large\bfseries\color{ink} O006 · 双缝一侧插入透明薄片的条纹平移}\quad{\small\color{dim}[难度 3 / 5 · 推导题]}\par
\noindent\textbf{题目}：双缝间距 $d$，屏距 $D\gg d$，用真空中波长为 $\lambda$ 的单色光照明。在一条缝前插入厚度为 $t$、折射率为 $n>1$ 的均匀薄片。忽略薄片引起的折射几何偏转并采用小角度近似，求干涉条纹在屏上平移的距离大小和对应的条纹数。
\par\noindent\textbf{解答}：薄片引入附加光程差 $\delta=(n-1)t$，与横向坐标引起的几何光程差 $dx/D$ 相抵。要使某条亮纹恢复原有级次，几何差变化量的绝对值应为 $\delta$，故移动距离 $|\Delta x|=D\delta/d$。原条纹间距 $\Delta x_0=\lambda D/d$，移动的条纹数为 $|\Delta x|/\Delta x_0=\delta/\lambda$。
\par\noindent\textbf{结论}：\(\displaystyle |\Delta x|=\frac{D(n-1)t}{d},\quad N=\frac{(n-1)t}{\lambda}\)\par
\vspace{1.0em}\par\noindent{\color{accent}\rule{\linewidth}{.35pt}}\vspace{.8em}
\section{近代物理}
\noindent{\large\bfseries\color{ink} Q001 · 光电效应截止电压}\quad{\small\color{dim}[难度 2 / 5 · 计算题]}\par
\noindent\textbf{题目}：某金属逸出功为 $2.0\,\mathrm{eV}$，照射光子能量为 $3.0\,\mathrm{eV}$。求光电子最大初动能及遏止电压的大小。
\par\noindent\textbf{解答}：爱因斯坦光电效应方程 $K_{\max}=E_\gamma-\phi=1.0\,\mathrm{eV}$。遏止电压 $U_s$ 满足 $eU_s=K_{\max}$，因 $1\,\mathrm{eV}=e\times1\,\mathrm V$，故 $U_s=1.0\,\mathrm V$。
\par\noindent\textbf{结论}：\(\displaystyle K_{\max}=1.0\,\mathrm{eV},\quad U_s=1.0\,\mathrm V\)\par
\vspace{1.0em}\par\noindent{\color{accent}\rule{\linewidth}{.35pt}}\vspace{.8em}
\noindent{\large\bfseries\color{ink} Q002 · 加速电子的德布罗意波长}\quad{\small\color{dim}[难度 3 / 5 · 计算题]}\par
\noindent\textbf{题目}：电子从静止出发，经 $150\,\mathrm V$ 电势差加速。忽略相对论修正，求其德布罗意波长。取 $h=6.626\times10^{-34}\,\mathrm{J\,s}$、$m_e=9.109\times10^{-31}\,\mathrm{kg}$、$e=1.602\times10^{-19}\,\mathrm C$。
\par\noindent\textbf{解答}：电子获得动能 $eU=p^2/(2m_e)$，所以 $p=\sqrt{2m_eeU}$，由德布罗意关系 $\lambda=h/p$。代入数值后为约 $1.00\times10^{-10}\,\mathrm m$，对应电子速度远低于光速，非相对论近似适用。
\par\noindent\textbf{结论}：\(\displaystyle \lambda\simeq0.100\,\mathrm{nm}\)\par
\vspace{1.0em}\par\noindent{\color{accent}\rule{\linewidth}{.35pt}}\vspace{.8em}
\noindent{\large\bfseries\color{ink} Q003 · 一维无限深势阱跃迁}\quad{\small\color{dim}[难度 3 / 5 · 推导题]}\par
\noindent\textbf{题目}：质量为 $m$ 的粒子处于宽度为 $L$ 的一维无限深势阱中，能级为 $E_n=n^2\pi^2\hbar^2/(2mL^2)$。求粒子由 $n=2$ 跃迁到 $n=1$ 时发射光子的频率。
\par\noindent\textbf{解答}：能量守恒给出 $h\nu=E_2-E_1=(4-1)\pi^2\hbar^2/(2mL^2)$。用 $h=2\pi\hbar$ 化简可得 $\nu=3\pi\hbar/(4mL^2)$。
\par\noindent\textbf{结论}：\(\displaystyle \nu=\frac{3\pi^2\hbar^2}{2mL^2h}=\frac{3\pi\hbar}{4mL^2}\)\par
\vspace{1.0em}\par\noindent{\color{accent}\rule{\linewidth}{.35pt}}\vspace{.8em}
\noindent{\large\bfseries\color{ink} Q004 · 放射性衰变与活度}\quad{\small\color{dim}[难度 2 / 5 · 计算题]}\par
\noindent\textbf{题目}：某放射性核素的半衰期为 $T_{1/2}$，初始核数为 $N_0$、初始活度为 $A_0$。经过两个半衰期后，求剩余核数、累积衰变核数和活度。假设无新核产生。
\par\noindent\textbf{解答}：衰变规律 $N(t)=N_0 2^{-t/T_{1/2}}$。当 $t=2T_{1/2}$，$N=N_0/4$，衰变核数 $N_0-N=3N_0/4$。由于活度 $A=\lambda N$ 且 $\lambda=\ln2/T_{1/2}$ 恒定，所以 $A/A_0=N/N_0=1/4$。
\par\noindent\textbf{结论}：\(\displaystyle N=N_0/4,\quad N_{\mathrm{衰变}}=3N_0/4,\quad A=A_0/4\)\par
\vspace{1.0em}\par\noindent{\color{accent}\rule{\linewidth}{.35pt}}\vspace{.8em}
\noindent{\large\bfseries\color{ink} Q005 · 高速不稳定粒子的存活比例}\quad{\small\color{dim}[难度 3 / 5 · 推导题]}\par
\noindent\textbf{题目}：某不稳定粒子的固有平均寿命为 $\tau_0$，相对实验室以恒速 $v=\beta c$（$0<\beta<1$）直线运动，其固有时衰变服从指数分布。求实验室中平均寿命、平均飞行距离，以及飞过长度 $L$ 后尚未衰变的比例。
\par\noindent\textbf{解答}：洛伦兹因子 $\gamma=(1-\beta^2)^{-1/2}$。时间膨胀给出实验室平均寿命 $\tau=\gamma\tau_0$，故平均飞行距离 $\bar\ell=v\tau=\beta c\gamma\tau_0$。飞越距离 $L$ 所需实验室时间为 $L/(\beta c)$，对应固有时间为 $L/(\beta c\gamma)$。指数存活律给出 $P=\exp[-L/(\beta c\gamma\tau_0)]$。
\par\noindent\textbf{结论}：\(\displaystyle \tau=\gamma\tau_0,\quad\bar\ell=\beta c\gamma\tau_0,\quad P(L)=e^{-L/(\beta c\gamma\tau_0)}\)\par
\vspace{1.0em}\par\noindent{\color{accent}\rule{\linewidth}{.35pt}}\vspace{.8em}
\noindent{\large\bfseries\color{ink} Q006 · 改变照射频率后的遏止电压}\quad{\small\color{dim}[难度 2 / 5 · 计算题]}\par
\noindent\textbf{题目}：某金属光电阈频为 $f_0$。用频率为 $f_1>f_0$ 的单色光照射时，遏止电压为 $U_1$。保持其他条件不变，将入射频率提高到 $f_2>f_1$，求新遏止电压及增量。电子电荷量的绝对值为 $e$。
\par\noindent\textbf{解答}：爱因斯坦光电方程 $eU_s=hf-hf_0$，所以 $U_1=h(f_1-f_0)/e$、$U_2=h(f_2-f_0)/e$。两式相减，金属逸出功抵消，$U_2-U_1=h(f_2-f_1)/e$。假设单光子光电效应、无接触电势修正。
\par\noindent\textbf{结论}：\(\displaystyle U_2=U_1+\frac{h(f_2-f_1)}e,\quad\Delta U=\frac{h(f_2-f_1)}e\)\par
\vspace{1.0em}\par\noindent{\color{accent}\rule{\linewidth}{.35pt}}\vspace{.8em}
\end{document}
